\clearpage
\VerdierPhysicalPageAnchor{0206}%
\begingroup\small\setlength{\parskip}{0pt}\setlength{\emergencystretch}{1.5em}
\setlength{\abovedisplayskip}{4pt plus 2pt minus 2pt}%
\setlength{\belowdisplayskip}{4pt plus 2pt minus 2pt}%

\VerdierSemanticAnchor{ch3d-2-4-5}
\noindent\textbf{Remark 2.4.5.}\quad
a) Let $\mathcal A$ be an abelian category all of whose objects are
isomorphic to quotients of projective objects. A complex of projective type
need not be $Q$-free on the left.

\VerdierSemanticAnchor{ch3d-2-4-5-example}
\emph{Example}: Let $\mathcal A$ be the category of
$\mathbb Z/4\mathbb Z$-modules. The complex
\[
\cdots\longrightarrow\mathbb Z/4\mathbb Z
\longrightarrow\mathbb Z/4\mathbb Z
\longrightarrow\mathbb Z/4\mathbb Z
\longrightarrow\mathbb Z/4\mathbb Z
\longrightarrow\mathbb Z/4\mathbb Z\longrightarrow\cdots
\]
in which every differential is multiplication by $2$ is acyclic and of
projective type. It is nevertheless not homotopic to zero. The same
complex gives an example of an acyclic complex of injective type that is
not homotopic to zero.

\VerdierSemanticAnchor{ch3d-2-4-5-b}
b) Even when the projective dimension of the objects of $\mathcal A$ is
not bounded, there may be complexes of projective type that are $Q$-free
on the left whose image in $D(\mathcal A)$ does not belong to
$D^{-}(\mathcal A)$. For example, a Cartan--Eilenberg projective complex
is an object that is $Q$-free on the left. More generally, every direct sum
of objects that are $Q$-free on the left is $Q$-free on the left.

\VerdierSemanticAnchor{ch3d-2-4-6}
\medskip\noindent\textbf{2.4.6.}\quad
We leave it to the reader to formulate the statements obtained by applying
the results of this section to the category opposite to the category
$\mathcal A$.
\endgroup

